% NAF 5166, batch 4 — Gallica views 61–80.
% Pass: Opus 5.5 (claude-opus-5-5), 2026-09-23 — first pass, unchecked against the views by a human.
%
% What the views are. Greyscale scans, portrait, about 4700 × 6430, one
% page per view; old leaves of several sizes hinged on stubs over larger
% modern guard leaves, as in batch 1. Where a leaf is folded over its guard
% the camera sometimes catches two pages in one view, the next page cut off
% at the edge (v. 63, 64, 67, 68): each page is transcribed once, at the
% view where it is whole. Every view was first looked at as a 1000-px
% thumbnail, the empty ones again at stretched contrast; each written page
% was cropped to the leaf and read as four full-width bands at 2400 px (12 %
% overlap), with tight crops at 1400–2400 px for every struck, interlined or
% doubtful word. Show-through was read mirrored where it bore on the order
% of the leaves (v. 61, 62, 69). Read only from the local mirror. \page{N}
% is always the view.
%
% What the twenty views hold.
%
%   v. 61     A small leaf, narrower than its guard, mounted in the lower two
%             thirds of the view, in verso position (the stubs are at the
%             right). One side of a draft letter, unsigned and undated,
%             « Vous [avez la bonté de, struck] m'engagez à reprendre une
%             tache que j'ai déja fort mal remplie… », to a « Monsieur »;
%             « je vous prie, Monsieur d'ageer, l'expression de ma
%             reconnoissance » is « d'agréer » with the r left out, not a
%             name (coordinator, from a tight crop, 2026-09-23); thanks for
%             kindness shown after an error, heavily corrected, breaking off
%             on a struck « que vous me témoignez » with no closing formula.
%             No number on the leaf. Its other side is written and shows
%             through mirrored; see « Where the pieces begin and end ».
%   v. 62–68  PIECE 2: a letter on leaves « 23 »–« 26 », one continuous
%             text: 23r (v. 62), 23v (v. 63), 24r (v. 64), 24v (v. 65), 25r
%             (v. 66), 25v (v. 67), 26r (v. 68). It begins at the head of
%             23r, « Je recommande à votre protection le mémoire No 1 portant
%             cette épigraphe », with a two-line Latin epigraph signed
%             « Baco novum org. », and ends two-thirds of the way down 26r:
%             « …d'agréer l'assurance des sentimens … d'admiration et de
%             respect de votre servante ». No name, no date, no place; the
%             addressee is « Monsieur » and is not named. The writer asks
%             the addressee's protection for her memoir before « la
%             classe » and its « commissaires »: its epigraph; the opinion
%             « M^r De la Grange » had given of the question; her theory
%             compared with « M^r Chladni »'s experiments and the one
%             discrepancy she admits (sounds from two particular integrals,
%             one of the form of the vibrating-string equation); the § 3 on
%             the words « appuyé » and « libre » in Euler's memoir; the
%             difference between nodal points and lines of lames and plates
%             and those of the string and the stretched surface (Euler,
%             Ricati); a passage of the « programme » quoted and to be
%             changed (Chladni, the cords of Sauveur); Sauveur's points of
%             rest; the § 6 on a soft body applied to the sides of a
%             rectangular plate; the last chapter on Euler's theory of the
%             ring against Chladni's experiments; her regret at not having
%             « l'avis de M^r De la Grange ». The page ends are sentence
%             breaks (« enfin de », « en montrant », « les conditions »,
%             « fort indépendant »): nothing is missing between the leaves.
%             A draft rather than a fair copy: many deletions and
%             interlinear corrections, but the flow is continuous.
%   v. 69     Verso of leaf « 26 », blank: mirrored show-through of 26r only
%             (the « 26 » shows reversed at the top left; a white fibre on
%             the guard beside it is a paper flaw). Not transcribed.
%   v. 70–79  Blank modern guard leaves, rectos and versos (checked at
%             contrast). Not transcribed.
%   v. 80     Inside of the lower board (pastedown, marbled edges), blank.
%             Not transcribed.
%   v. 81 (batch 5): the back board of the binding (marbled paper, spine
%             with the shelf label); not transcribed, and batch 5 has no
%             file.
%
%   Where the pieces begin and end, as far as this batch shows.
%   - v. 61 opens at the head of its page with a capital « Vous », with no
%     sentence carried over from before. But the leaf sits in verso
%     position, and its other side is written: mirrored through the paper,
%     some ten lines, of which can be made out « à la place de y et z, je
%     sais bien que ce n'est pas », « M^r De la Grange », « frappée à la vue
%     de son équation », « je n'ai pû résister à la tentation de vous la
%     communiquer », « en protestant au reste de ma parfaite soumission »,
%     « aux oracles », and a last line ending « géomètre. ». That side is
%     presumably view 60 (batch 3), which this pass has not seen. Whether it
%     and v. 61 are one draft or two, and which comes first, cannot be told
%     from here; the v. 61 text is transcribed as a piece of its own, not as
%     a continuation.
%   - v. 61 and v. 62 are different leaves: v. 61's text breaks off
%     unfinished; v. 62 begins a new letter.
%   - v. 80 is the lower board; nothing runs on into v. 81, the back board.
%   - The letter of v. 62–68 is complete in this batch: first words at the
%     head of 23r, subscription on 26r.
%
%   The hand. One sloping French cursive throughout v. 61–68, apparently the
%   same on the small leaf of v. 61 and on leaves 23–26: long s (« ſens »,
%   « diſſimuler », « ſuccès », « ſentence »), transcribed as s; words tied
%   by long hair-lines at the baseline, so that a baseline stroke is not a
%   deletion (a deletion runs through the body of the letters); t-bars long
%   enough to cross the next letter; a looped capital M in « Monsieur »,
%   « Mais », « M^r ». The writer's grammar is feminine (« disposée »,
%   « fondée », « contente », « trompée », « égarée », « sûre », « votre
%   servante »). It is not the hand of the memoir of v. 17–19 (batch 1):
%   no long s there, a quicker, heavier ductus. The pass does not say whose
%   hand this is: no name is written on these leaves, and the wrapper's
%   mention of « notes … de Mlle Sophie Germain » (v. 15) is the wrapper's.
%   Spelling as written: « avoit », « auroit », « seroit », « connoitre »,
%   « reconnoissance », « parceque », « parconséquent », « indépendans »,
%   « sentimens », « d'avantage », « dédomagement », « aprouvé »,
%   « employe », « fesoient », « ennuieux », « Ricati », « isdem ».
%
%   Notation. No formulae in these views. « § » is written as two linked s
%   and set §. Underlinings are set \emph{}.
%
%   Numbers on the leaves, and which kind.
%   - « 23 » (v. 62), « 24 » (v. 64), « 25 » (v. 66, the 5 like an S),
%     « 26 » (v. 68): ink, top right of the rectos, one cursive series, the
%     same kind of figure as batch 1's « 2 » (v. 19), not the posed « 1 »
%     of v. 17. They run consecutively across one letter, and « 26 » is the
%     last of the « 26 Feuillets » of the 1888 certificate (v. 13), so they
%     behave as the certificate's series; being ink, and following batch 1
%     and the NAF 5161 practice, no \folio{} is written. No pencilled
%     figure was seen on these leaves.
%   - No number on the leaf of v. 61, nor on the versos (v. 63, 65, 67, 69).
%   - In the text: « N^o 1 » (the memoir entered, v. 62), « § 3 » (v. 64),
%     « § 6 » (v. 66) — the writer's references to her own memoir's
%     paragraphs, not leaf numbers.
%   - Stamps « BIBLIOTHÈQUE NATIONALE / R.F. / MANUSCRITS » on v. 62, 66, 68.
%
%   Apparatus. Five \ill{}, all inside \struck{} (v. 61 ×1, 63 ×1, 66 ×1,
%   67 ×2); ten \uncertain{}. One \add{} (a closing guillemet, v. 65).

% Coordinator's reconciliation (2026-09-23): the other side of the v. 61
%   leaf is v. 60 (batch 3), leaf 22 recto, the end of the draft letter of
%   v. 58–60; the phrases read here in show-through are in batch 3's
%   transcription. The numbers 23–26 continue the library's series 1–22
%   (see batch 1's reconciliation); no \folio{} is written.

\documentclass[11pt,a4paper]{article}
% The preamble shared by every transcription of the mathematicians' manuscripts in Gallica.
%
% It rests on one idea: **uncertainty compounds, it does not resolve.** A
% transcription that smooths over a doubtful word has destroyed the only thing
% it was made for — a reader must be able to tell, at every point, what was on
% the page and what was supplied. The apparatus macros below are therefore the
% critical apparatus, not decoration.
%
% The same commands are recognised by `scripts/render.mjs`, which produces the
% site's reading view, and by `scripts/tei.mjs`. A macro added here must be
% added there too, or rendering fails loudly — which is the wanted behaviour.
%
% Comments here are English, like the rest of the repository. The documents
% this preamble sets are French, and that is deliberate: see the skills.

\ProvidesPackage{germain}[2026/09/15 Transcriptions of mathematicians' manuscripts in Gallica]

\usepackage{amsmath,amssymb,amsthm}
\usepackage{mathrsfs}
% Kept from the parent preamble so the renderer's diagram support stays
% available; these manuscripts draw no commutative diagrams, and nothing here uses it.
\usepackage{tikz-cd}
\usepackage[english,french]{babel}
\usepackage{fontspec}

% --- Greek ------------------------------------------------------------------
% The text face stays fontspec's default, Latin Modern, which has no Greek:
% XeTeX drops every glyph it lacks with a line in the log and nothing on the
% page, and the Greek words the manuscripts quote vanished from the PDFs. So
% the Greek and Greek Extended blocks get a XeTeX character class of their own,
% and entering or leaving a run of it switches the family to CMU Serif — the
% same Computer Modern design, with polytonic Greek — and back. Only the
% family changes: a Greek word in \textit or \textbf keeps its shape and series.
% The transitions to and from every other class carry the tokens that class
% has towards an ordinary letter, so French spacing before « ; » or inside
% « » still holds after a Greek word. No grouping, so no brace or \endgroup
% can come out unbalanced; maths is left alone, since XeTeX inserts nothing
% there. Kept identical in `exercices.sty`.
\makeatletter
\newfontfamily\germainGreekfont{cmun}[
  Extension=.otf, NFSSFamily=germaingreek,
  UprightFont=*rm, ItalicFont=*ti, SlantedFont=*sl,
  BoldFont=*bx, BoldItalicFont=*bi, BoldSlantedFont=*bl]
\newXeTeXintercharclass\germain@greek
\def\germain@greekrange#1#2{%
  \@tempcnta=#1\relax
  \loop
    \XeTeXcharclass\@tempcnta=\germain@greek
    \ifnum\@tempcnta<#2\advance\@tempcnta\@ne\repeat}
\germain@greekrange{"0370}{"03FF}
\germain@greekrange{"1F00}{"1FFF}
\def\germain@greekprev{\rmdefault}
\def\germain@greekin{%
  \edef\germain@greekprev{\f@family}\fontfamily{germaingreek}\selectfont}
\def\germain@greekout{\fontfamily{\germain@greekprev}\selectfont}
\def\germain@greekjoin{%
  \XeTeXinterchartokenstate=\@ne
  \@tempcnta=\z@
  \loop
    \ifnum\@tempcnta=\germain@greek\else
      \XeTeXinterchartoks\germain@greek\@tempcnta=\expandafter{%
        \expandafter\germain@greekout\the\XeTeXinterchartoks\z@\@tempcnta}%
      \XeTeXinterchartoks\@tempcnta\germain@greek=\expandafter{%
        \the\XeTeXinterchartoks\@tempcnta\z@\germain@greekin}%
    \fi
    \ifnum\@tempcnta<4095\advance\@tempcnta\@ne\repeat}
% babel-french sets its own classes, and saves and restores the interchar
% state, each time a language is selected: join again after it does.
\addto\extrasfrench{\germain@greekjoin}
\addto\extrasenglish{\germain@greekjoin}
\AtBeginDocument{\germain@greekjoin}
\makeatother

\usepackage{geometry}
\geometry{a4paper,margin=25mm,marginparwidth=28mm}
\usepackage{marginnote}
\usepackage{xcolor}
\usepackage[normalem]{ulem}
\setlength{\emergencystretch}{3em}
\usepackage{hyperref}
\hypersetup{colorlinks=true,linkcolor=black,urlcolor=black,pdfborder={0 0 0}}

% --- Batch metadata ---------------------------------------------------------
% Taken as they stand by the rendering script, which shows them at the head of
% the reading view. They fix what the file claims to cover. The macro names are
% the parent project's, so that the scripts stay shared; \folder here means the
% volume's slug — `fr-9115` — and \pages counts Gallica views.

\newcommand{\folder}[1]{\gdef\grothFolder{#1}}
\newcommand{\batch}[1]{\gdef\grothBatch{#1}}
\newcommand{\pages}[2]{\gdef\grothFirst{#1}\gdef\grothLast{#2}}
\newcommand{\foldertitle}[1]{\gdef\grothFoldertitle{#1}}

% The holder's own shelfmark, printed as they write it: « Français 9115 ».
\newcommand{\shelfmark}[1]{\gdef\grothShelfmark{#1}}

% The dating, copied verbatim from the holder's catalogue — never inferred
% here. « XIXe siècle » is the BnF's; a pass that dates a leaf from its content
% says so in a \note{}, as its own inference.
\newcommand{\dating}[1]{\gdef\grothDating{#1}}

% The status notice, printed in the title block and repeated as a diagonal
% watermark on every page of the PDF. The manuscripts are in the public
% domain; what this line declares is not a right but a quality — an
% unverified first pass by a machine. The skills write the line; a file
% without it gets no watermark, so leaving it out is visible in the output.
\newcommand{\watermark}[1]{\gdef\grothWatermark{#1}}

\gdef\grothFolder{?}\gdef\grothBatch{}\gdef\grothFirst{?}\gdef\grothLast{?}%
\gdef\grothFoldertitle{}\gdef\grothShelfmark{}\gdef\grothDating{}\gdef\grothWatermark{}

% --- The résumé -------------------------------------------------------------
\newenvironment{resume}
  {\small\setlength{\parskip}{0.45em}}
  {\par\medskip\hrule\medskip}

% --- Keywords ---------------------------------------------------------------
% In English, closing the modernised reading's résumé: the modern vocabulary
% under which to search for what the leaves construct. The manifest extracts
% them to tag the volume — this line is the single source of the tags.
\newcommand{\keywords}[1]{\par\smallskip\noindent
  {\footnotesize\textsc{Keywords} --- \textit{#1}}\par}

% --- The critical apparatus -------------------------------------------------

% The Gallica view number, in the margin. This is the anchor that lets a
% disagreement over a formula be settled by looking at *one* image rather than
% a volume of 750. The site uses it to turn the facsimile as the transcription
% is scrolled. A view is one scanned image: usually one side of a leaf.
\newcommand{\page}[1]{%
  \marginnote{\footnotesize\color{gray!70}\textsf{v.\,#1}}%
  \label{page:#1}\ignorespaces}

% The BnF's pencilled foliation, where the leaf carries it and a pass can read
% it: « 198r ». This is what the literature cites — « ff. 198r–208v » — and
% the one line that turns a citation into a view. Recorded, never inferred:
% a folio a pass cannot read is not written.
\newcommand{\folio}[1]{%
  \marginnote{\footnotesize\color{gray!50}\textsf{f.\,#1}}\ignorespaces}

% The modernised reading's coarser counterpart to \page{}: which views a
% section is drawn from.
\newcommand{\pagerange}[2]{%
  \marginnote{\footnotesize\color{gray!70}\textsf{v.\,#1--#2}}%
  \ignorespaces}

% Illegible: never guessed.
\newcommand{\ill}{\textcolor{red!60!black}{[\,\ldots\,]}}

% A doubtful reading: the word is given, and flagged as uncertain.
\newcommand{\uncertain}[1]{\uline{#1}}

% An editorial addition — a missing letter, an implied word.
\newcommand{\add}[1]{\textcolor{blue!50!black}{[#1]}}

% Struck out by the writer. What was crossed out is part of the document.
\newcommand{\struck}[1]{\sout{#1}}

% The transcriber's note, on the state of the page or the mathematics.
\newcommand{\note}[1]{\par\smallskip{\small\color{gray!60!black}\textit{[#1]}}\par\smallskip}

% A marginal note of hers, distinct from the body of the text.
\newcommand{\marginal}[1]{\par\smallskip\noindent\hspace{1em}%
  {\small\color{gray!60!black}#1}\par\smallskip}

% --- Title ------------------------------------------------------------------
% The title block is French, like the documents themselves.

\AddToHook{shipout/background}{%
  \ifx\grothWatermark\empty\else
    \begin{tikzpicture}[remember picture, overlay]
      \node[rotate=52, scale=3.4, align=center, text=gray!18,
            font=\sffamily\bfseries]
        at (current page.center) {\grothWatermark};
    \end{tikzpicture}%
  \fi}

\AtBeginDocument{%
  \begin{center}
    {\large\bfseries Manuscrits de mathématiciens (Gallica) ---
      \ifx\grothShelfmark\empty\grothFolder\else\grothShelfmark\fi}\\[2pt]
    {\itshape\grothFoldertitle}\\[4pt]
    {\small vues \grothFirst--\grothLast
      \ifx\grothBatch\empty\else\ (lot \grothBatch)\fi}\\[2pt]
    \ifx\grothDating\empty\else
      {\small\color{gray!60!black}Datation du catalogue : \grothDating}\\[2pt]
    \fi
    \ifx\grothWatermark\empty\else
      {\small\bfseries\color{red!45!gray}\grothWatermark}\\[2pt]
    \fi
    {\footnotesize\color{gray!60!black}Transcription établie d'après les images IIIF de Gallica.
     Source gallica.bnf.fr / Bibliothèque nationale de France.\\
     Les lectures incertaines sont soulignées, les illisibles notées [\,\ldots\,],
     les ajouts entre crochets ; \textsf{v.} numérote les vues de Gallica,
     \textsf{f.} la foliotation au crayon de la BnF.}
  \end{center}
  \vspace{1em}}

\endinput


\folder{naf-5166}
\shelfmark{NAF 5166}
\batch{4}
\pages{61}{80}
\dating{1701-1900}
\watermark{Lecture automatique\\non vérifiée}
\foldertitle{Papiers du géomètre J.-L. LAGRANGE (1813).}

\begin{document}

\page{61}
\note{Petit feuillet de papier ancien, plus étroit que le feuillet de garde moderne sur lequel il est monté, et qui n'en occupe que les deux tiers inférieurs. Aucun numéro, ni à l'encre ni au crayon, ne se lit sur le feuillet ou sur la garde. Le texte est un brouillon de lettre, d'une main cursive penchée, autre que celle du mémoire des vues 17--19, qui emploie le \emph{s} long (rendu par \emph{s}) ; les accords sont au féminin (« disposée », « fondée », « contente »). Ni date, ni en-tête, ni signature. L'autre face du feuillet est écrite : son texte transparaît, à l'envers, dans le bas de la page, mais cette face n'est pas une vue de ce lot ; voir l'en-tête.}

Vous \struck{avez la bonté de} m'engage\uncertain{z} à reprendre une tache que j'ai déja fort mal remplie ; je m'y sens peu disposée et il faut convenir que puisque j'ai pû me tromper si grossierement \struck{il} je ne suis que trop fondée à me méfier de mes forces, cependant j'aime assez la verité pour etre encore contente de la connoitre lors même qu'elle ne m'est pas favorable \struck{et} j'ai \struck{bien} d'ailleurs à me féliciter d'avoir \struck{une} obtenu, dans cette occasion, une nouvelle preuve de l'interet que vous daignez m'accorder, c'est \struck{bien sincerement} que je vous assure bien sincerement que \struck{la} votre bienveillance est pour moi un dédomagement \struck{d'un} \struck{auquel j'attache le plus} grand prix, et que j'apprécie, comme je le dois le soin que vous prenez de me dissimuler le ridicule attaché \struck{à la médiocrité} au mauvais succès, je vous prie, Monsieur d'ageer\note{« d'ageer » : ainsi sur la feuille, lettre à lettre, pour « d'agréer » ; l'écrivante a omis le \emph{r}. Ce n'est pas le nom du destinataire, qui n'est pas nommé.}, l'expression de ma reconnoissance et de croire que \struck{\ill{}} le respect et l'admiration que j'ai toujours eu pour vous donnent un prix infini à \uncertain{tant} de \struck{à la} bontés \struck{que vous me témoignez}
\note{« m'engagez » : la fin du mot est surchargée et barrée ; on lit « m'engage » suivi d'un \emph{z} ou d'un \emph{r} corrigé. Au-dessus de « reprendre », un grand \emph{D} bouclé isolé, dont le rôle ne se voit pas. « obtenu » : la fin du mot est repassée d'un trait. « nouvelle » est écrit dans l'interligne au-dessus de « preuve ». « d'ailleurs » est écrit dans l'interligne au-dessus du « bien » biffé. « bien sincerement » : biffé dans la ligne après « c'est », puis récrit dans la marge gauche, devant « que votre bienveillance ». « auquel j'attache le plus » est écrit dans l'interligne au-dessus de « d'un grand prix » et biffé ensuite ; « d'un » est biffé aussi, de sorte que la phrase finale lit « un dédomagement grand prix » sans correction complète. « à la médiocrité » : biffé d'un trait, lecture sûre pour « à la », probable pour « médiocrité ». « Monsieur \emph{d'ageer} » : le nom se lit \emph{d}, apostrophe (ou accent), \emph{a}, \emph{g}, \emph{e}, \emph{e} ou \emph{c}, \emph{r} ; il n'est pas identifié ici. Après « que », deux signes brefs biffés, le premier semblable à un \emph{k}, le second à une esperluette ; illisibles. « à tant de » est écrit dans l'interligne au-dessus de « à la bontés », le mot « tant » peut-être barré lui aussi ; « bontés » : le \emph{s} est ajouté sur le point final. La lettre s'arrête sur « témoignez », biffé, sans point ni formule.}

\page{62}
\note{Feuillet de papier ancien, monté sur onglet, le texte sur les trois cinquièmes supérieurs. En haut à droite, à l'encre, « 23 », d'un chiffre cursif ; aucun folio n'est porté, voir l'en-tête. Le cachet rond « BIBLIOTHÈQUE NATIONALE / R.F. / MANUSCRITS » est posé à droite, sur « manifestée ». Même main que la vue 61 (\emph{s} long, rendu par \emph{s}). Dans le haut de la page transparaît, à l'envers, le début du verso, qui est la vue 63. Ni en-tête, ni date, ni nom de destinataire sur cette page.}

Je recommande à votre protection le mémoire N\textsuperscript{o} 1 portant cette épigraphe :

\begin{quote}
« Sed longe maximum progressibus scientiarum, et novis pensis ac provinciis in isdem suscipiendis « obstaculum deprehenditur in desperatione hominum et suppositione impossibilis. \emph{Baco novum org.}
\end{quote}
\note{L'épigraphe tient deux lignes, chacune ouverte par des guillemets en marge ; « Baco novum org. » est souligné. « N\textsuperscript{o} 1 » : un \emph{N} à \emph{o} suscrit, puis un chiffre bref, lu « 1 ». « isdem » : ainsi, un seul \emph{i}.}

Si \struck{j'eusse} j'\uncertain{en} avois trouvé l'occasion je vous aurois consulté avant de l'adopter, car elle a un air de \uncertain{venterie} qui ne me convient guère à moi qui ai tant de raisons de me défier de mes forces et qui ne vois pas en effet de plus forte objection contre ma théorie que l'improbabilité d'avoir rencontré juste. Cependant je crains l'influence de l'opinion qu'avoit manifestée M\textsuperscript{r} De la Grange ; la question n'a sans doute été abandonnée que parceque ce grand géomètre l'avoit jugée difficile ; peut-être le même préjugé me fera-t-il condamner sans un examen réfléchi et c'est ce qui m'a engagé à placer à la tête de mon mémoire une sentence qui semble faite exprès pour la disposition des esprits à cet égard. Au reste, je compte d'avantage sur votre appui que sur l'effet de la pensée philosophique de Bacon et quoique le sujet du mémoire ne me laisse pas l'espérance de vous avoir pour juge, je me recommande à vous, pour faire en sorte que les commissaires prennent la peine d'entendre mon long et ennuieux travail.
\note{« j'en avois » est écrit dans l'interligne au-dessus de « j'eusse », biffé ; le « en » est peut-être barré lui aussi. « venterie » : lettres \emph{v}, \emph{e}, \emph{n}, \emph{t}, puis une voyelle ronde, \emph{r}, \emph{i}, \emph{e} ; lu avec doute. « M\textsuperscript{r} De la Grange » : en deux mots, « De la » en fin de ligne et « Grange » au début de la suivante. « sentence » : le \emph{n} du milieu est repassé. « d'avantage » : ainsi. La page s'arrête sur « travail. » ; la lettre continue au verso (vue 63).}

\page{63}
\note{Verso du feuillet « 23 » (vue 62), à gauche de la vue ; aucun numéro n'y est écrit. À droite, le recto du feuillet suivant, coté « 24 », paraît en partie, coupé par le bord de l'image : il est transcrit à la vue 64, où on le voit entier. La lettre continue ici sans interruption depuis « travail. » (vue 62).}

Je crois être sûre de la théorie qui en est l'objet, je l'ai examinée un grand nombre de fois, je l'ai comparée aux résultats des expériences de M\textsuperscript{r} Chladni et je n'ai pas dissimulé, au reste, l'objection que l'on peut tirer d'une différence entre les rapports déduits de cette théorie et les rapports conclus de l'expérience. Mais cette différence ne se trouve que dans certains cas ; ce sont ceux qui appartiennent à l'intégrale particulière qui donne lieu à une équation de la forme de celle aux cordes vibrantes ; elle est d'ailleurs peu considérable, et ne se fait sentir que dans la comparaison des sons qui accompagnent les figures renfermées dans cette intégrale, aux sons qui accompagnent les figures qui résultent d'une autre intégrale particulière. En effet les sons donnés par chacune de ces intégrales, pris dans les mêmes circonstances, donnent, en les comparant entre eux, \struck{absolument} les \struck{mêmes} rapports \struck{que} \struck{absolument} \struck{entièrement} conformes à l'expérience.
\note{« de cette théorie » : « cette » est surchargé d'un trait épais, peut-être biffé ; il est gardé. « renfermées » : écrit « renfermés », l'accord n'est pas fait. Le premier « absolument » biffé est en fin de ligne, après « entre eux, » ; le second, avec « entièrement », dans la ligne suivante. « mêmes » : biffé, lecture probable.}

A l'égard des figures j'en ai expliqué un grand nombre \struck{de la} d'une manière \struck{la plus} que je crois satisfaisante : ainsi je \struck{crois} pense \struck{la théorie ap} que ma théorie est \struck{\ill{}} appuyée sur des preuves suffisantes, et qu'elle est même plus avancée, sous le rapport de la comparaison à l'expérience que celle des surface tendues sur laquelle personne n'élève de doute.
\note{« d'une », « que je crois » et « pense » sont écrits dans l'interligne, au-dessus de « de la », « la plus » et « crois », biffés. Avant « appuyée », un mot biffé, commençant peut-être par \emph{r}, illisible, suivi d'un groupe bref qui ressemble à « dûe », non barré, dont on ne voit pas s'il appartient au mot biffé. « des surface tendues » : ainsi. Le verso s'arrête sur « doute. » au pied de la page ; la lettre continue au feuillet « 24 » (vue 64).}

\page{64}
\note{Recto du feuillet coté « 24 », à l'encre, en haut à droite, d'un chiffre cursif semblable au « 23 » de la vue 62 ; aucun folio n'est porté, voir l'en-tête. À gauche paraît en partie le verso du feuillet « 23 », déjà transcrit à la vue 63. Même main.}

Mais lors même que je me serois égarée dans l'objet \uncertain{principal} de mes recherches, il resteroit encore dans mon mémoire des § qui ne seroient peut-être pas indignes de l'attention de la classe. Par exemple, \struck{le} § 3 \struck{dans lequel j'ai \uncertain{fait examen}} de la valeur analytique des mots \emph{appuyé} et \emph{libre} qui ont été employés jusqu'ici d'une manière tellement vague qu'il en est résulté, même dans le mémoire d'Euler, une grande confusion d'idée. Les remarques que j'ai faites dans le même chapitre sur la différence entre les points et \struck{les} lignes de repos qui se forment pendant le mouvement\struck{s} de la lame\struck{s} et de la plaque\struck{s} élastiques vibrantes et les points et lignes de repos que l'on observe dans le\struck{s} mouvement de la corde et de la surface tendues : différence qui tient à l'ordre des équations différentielles de ces corps sonores, et qui, par rapport à la lame, résulte à la fois, de la théorie d'Euler examinée de plus près qu'on ne l'avoit encore fait, du travail de Ricati qui en calculant les valeurs que cette théorie doit fournir a pourtant négligé de les comparer, enfin de
\note{« principal » : le mot se lit \emph{princip}, puis une fin brève (\emph{-eus} ou \emph{-al}) ; lu avec doute. « des § » : le signe de paragraphe, deux \emph{s} liés. « le § 3 dans lequel j'ai fait examen » : « le » est barré nettement ; « dans lequel j'ai fait examen » est traversé d'un trait continu à mi-hauteur des lettres, que l'on prend pour une biffure, et les deux derniers mots se lisent avec doute ; la phrase qui reste, « § 3 de la valeur analytique des mots », cite le titre d'un paragraphe. « appuyé » et « libre » sont soulignés. « mouvement », « lame », « plaque », « le » : un \emph{s} final (ou un \emph{x}) est barré à chacun, le pluriel ayant été ramené au singulier ; « élastiques vibrantes » est resté au pluriel. « Ricati » : ainsi. La page s'arrête sur « enfin de » ; la phrase continue au verso (vue 65).}

\page{65}
\note{Verso du feuillet « 24 » ; aucun numéro n'y est écrit. Même main. La phrase continue depuis « enfin de » (vue 64).}

de l'expérience qui, lorsqu'on employe des lames suffisamment longues, présente entre chaque ligne de repos des intervalles sensiblement inégaux et de longueur respective parfaitement conforme à ce qu'indique la théorie d'Euler ; quoique M\textsuperscript{r} Chladni, qui sans doute s'est \struck{un} servi de lames trop courtes, n'ait pas fait mention de cette inégalité des parties vibrantes.
\note{Le mot biffé avant « servi », en tête de ligne, se lit « un » ou « en ».}

\struck{Il résulte de mes remarques} Je crois donc, que lors même que je me serois trompée sur la question des surfaces, la classe devroit toujours, dans \struck{le} cas qu'elle soit encore remise, changer le passage de son programme dans lequel on lit « il a découvert (\struck{il s'agit des expériences de} Chladni) et rendu « perceptible, d'une manière très ingénieuse dans ces lames, des « \emph{\uncertain{nappes} vibrantes} analogues aux \emph{ondes} des cordes de Sauveur, « et des courbes \emph{d'équilibres} ou de repos auxquels correspondent « les \emph{nœuds} ou \emph{points de repos} des mêmes cordes. \add{»}
\note{« Je crois donc » est écrit dans l'interligne au-dessus de « Il résulte de mes remarques », biffé. « dans le cas qu'elle soit encore remise » est écrit dans la marge droite, sur trois lignes serrées, à la hauteur de « toujours, » ; « le » y est biffé, et au-dessus de « cas » deux lettres brèves, « au » ou « ou », avec un signe d'insertion, dont la place dans la phrase n'est pas sûre. La citation du programme est ouverte, comme l'épigraphe de la vue 62, par des guillemets en tête de chaque ligne ; elle ne se ferme que par un point après « cordes » ; le guillemet fermant est ajouté. Dans la parenthèse, « il s'agit des expériences de » est biffé, « Chladni » gardé. Les mots en italique sont soulignés sur la page. « nappes » : le mot souligné se lit \emph{n}, une voyelle, deux lettres à hampe descendante, \emph{e}, \emph{s} ; lu avec doute. « d'équilibres » : ainsi, au pluriel.}

En effet, le caractère des points de repos observés par Sauveur est que chacun d'eux soit une véritable limite analytique des parties qu'il sépare, en sorte que, si, par exemple, on partageoit effectivement en deux parties égales entr'elles la corde qui vibre en montrant
\note{La page s'arrête sur « en montrant » au pied du verso ; la phrase continue à la vue 66.}

\page{66}
\note{Recto du feuillet coté « 25 », à l'encre, en haut à droite, d'un chiffre cursif comme « 23 » et « 24 » (le 5 a la forme d'un \emph{S}) ; aucun folio n'est porté, voir l'en-tête. Cachet « BIBLIOTHÈQUE NATIONALE / R.F. / MANUSCRITS » à droite. Même main. La phrase continue depuis « en montrant » (vue 65).}

dans son milieu un seul nœud de vibration, les deux parties isolées, et soumises d'ailleurs aux mêmes conditions, continueroient de se ployer aux mêmes courbures et de donner les mêmes sons que lorsqu'elles fesoient parties de la corde entière.

Tandis qu'au contraire, lorsqu'il s'agit de la lame, il y a des cas, et c'est le plus grand nombre, où les parties \struck{\ill{}} séparées par des points de repos, si elles étoient isolées les unes des autres, ne pourroient jamais se ployer aux mêmes courbures, ni parconséquent rendre les mêmes sons que lorsqu'elles font parties de la lame entière.
\note{« continueroient » : quelques lettres du milieu sont repassées. Après « les parties », en fin de ligne, un mot biffé commençant par \emph{s} et portant deux hampes descendantes, peut-être un premier « séparées » abandonné ; illisible.}

Le § 6, qui est un véritable hors d'\uncertain{œuvre} au sujet du mémoire contient l'exposition d'un fait singulier dont j'ai eu l'honneur de vous parler, il consiste en ce que, au moyen de l'application d'un corps mou aux côtés d'une plaque rectangulaire élastique, on peut tirer d'une portion aussi petite que l'on veut, le même son que d'une portion de longueur donnée prise sur la même plaque. C'est-à-dire qu'à quelque distance de l'origine que se fasse la substitution de la matière non élastique, à une portion de la surface capable de faire équilibre à la substance ajoutée, les conditions
\note{« hors d'œuvre » : le mot s'écrit \emph{d'}, puis \emph{a} ou \emph{œ}, \emph{u}, \emph{v}, \emph{r}, \emph{e} ; lu avec doute. « § 6 » : le signe de paragraphe, puis un 6 net. La page s'arrête sur « les conditions » ; la phrase continue au verso (vue 67).}

\page{67}
\note{Verso du feuillet « 25 », à gauche de la vue ; aucun numéro n'y est écrit. À droite, le recto du feuillet suivant, coté « 26 », paraît en partie, coupé par le bord de l'image : il est transcrit à la vue 68. Même main. La phrase continue depuis « les conditions » (vue 66).}

des extrémités sont toujours satisfaites. On voit la raison de ce phénomène en remarquant que les termes qui doivent être nuls aux limites, sont l'un et l'autre multipliés par un facteur dépendant de la rigidité naturelle.

Enfin le dernier chapitre de mon mémoire contient \struck{une} l'explication des différences qui se trouvent entre les résultats de la théorie de l'anneau donnée par Euler et ceux des expériences de M\textsuperscript{r} Chladni : ainsi voila \struck{trois} plusieurs points parfaitement indépendans de la théorie des surfaces et qui sont appuyés tant sur la théorie d'Euler que sur des expériences faciles à répéter.
\note{« plusieurs » est écrit dans l'interligne au-dessus de « trois », biffé. Sous « multipliés », un petit signe en forme de chevron, sans appel visible.}

Malgré toutes les raisons que je crois voir en faveur de \struck{mon travail} mes idées j'ai si peu de confiance dans mon jugement que je doute encore de leur valeur. \struck{\ill{} \ill{}} Combien je regrette l'avis de M\textsuperscript{r} De la Grange ! s'il m'eut condamnée il auroit au moins remarqué ce qui a indépendamment de la théorie principale, peut mériter l'attention, s'il m'eut aprouvé j'\uncertain{aurais} été sûre d'avoir raison et peut-être mon travail tout imparfait qu'il est, lui auroit-il fourni l'occasion de voir de plus haut une théorie à laquelle je prends un intérêt fort indépendant
\note{« mes idées » est écrit dans l'interligne au-dessus de « mon travail », biffé (« mon » en fin de ligne, « travail » au début de la suivante). Avant « Combien », deux mots biffés d'un trait, le premier commençant par une majuscule bouclée, le second par \emph{r} ou \emph{reg} ; illisibles. « De la Grange » : en deux mots, comme à la vue 62 ; le point d'exclamation est net. « aurais » : la voyelle de la dernière syllabe se lit plutôt \emph{a} ; ailleurs la main écrit « auroit ». La page s'arrête sur « indépendant » au pied du verso ; la phrase continue au feuillet « 26 » (vue 68).}

\page{68}
\note{Recto du feuillet coté « 26 », à l'encre, en haut à droite, du même chiffre cursif que « 23 »--« 25 » ; aucun folio n'est porté, voir l'en-tête. À gauche paraît en partie le verso du feuillet « 25 », déjà transcrit à la vue 67. Même main. La phrase continue depuis « fort indépendant » (vue 67). Le texte occupe la moitié supérieure de la page ; le reste est blanc. Le cachet « BIBLIOTHÈQUE NATIONALE / R.F. / MANUSCRITS » est posé sur l'avant-dernière ligne.}

des calculs de mon amour propre. D'ailleurs l'idée de la difficulté de la question, qui empêchera peut-être que l'on ne se donne la peine d'examiner un \struck{travail} mémoire condamné d'avance, n'auroit pû alors \struck{m'inspirer} me donner aucune crainte. De pareils regrets sont aussi naturels qu'ils sont superflus, ils me rendent votre protection plus nécessaire et je la réclame avec la confiance que doit m'inspirer l'intérêt dont vous m'avez toujours honoré
\note{« mémoire » et « me donner » sont écrits dans l'interligne au-dessus de « travail » et de « m'inspirer », biffés. « d'examiner » : les premières lettres sont surchargées. Pas de point après « honoré ».}

Vous me pardonnerez peut-être la longueur de mon plaidoyer qui mérite bien que je vous en fasse excuse en vous priant, Monsieur, d'agréer l'assurance des sentimens de \struck{respect et} d'admiration et de respect de votre servante
\note{« et de » avant le second « respect » est sous le cachet ; un trait passe sur « de », qui peut être une biffure ou le trait de liaison de la main ; gardé. La lettre s'arrête sur « de votre servante », sans nom, sans date et sans lieu ; aucun nom de destinataire n'est écrit sur les feuillets « 23 » à « 26 ».}

\end{document}
